\documentclass[11pt, a4paper]{article}

% ── Packages ──────────────────────────────────────────────────────────
\usepackage[utf8]{inputenc}
\usepackage[T1]{fontenc}
\usepackage{lmodern}
\usepackage[margin=1in]{geometry}
\usepackage{amsmath, amssymb, amsthm}
\usepackage{booktabs}
\usepackage{longtable}
\usepackage{graphicx}
\usepackage{hyperref}
\usepackage{xcolor}
\usepackage{listings}
\usepackage{tcolorbox}
\usepackage{polya}

\hypersetup{
  colorlinks=true,
  linkcolor=polyablue,
  urlcolor=polyablue,
  citecolor=polyablue,
}

% ── Code listing style ────────────────────────────────────────────────
\lstdefinestyle{polya-python}{
  language=Python,
  basicstyle=\ttfamily\small,
  keywordstyle=\color{polyablue}\bfseries,
  commentstyle=\color{polyagray}\itshape,
  stringstyle=\color{polyagreen},
  numbers=left,
  numberstyle=\tiny\color{polyagray},
  numbersep=8pt,
  frame=single,
  framerule=0.4pt,
  rulecolor=\color{polyagray},
  breaklines=true,
  showstringspaces=false,
  tabsize=4,
}
\lstset{style=polya-python}
\title{Trip Itinerary}
\author{uber-polya}
\date{2026-02-26}

\begin{document}

\maketitle
\tableofcontents
\newpage

% ── Phase A: Problem Formulation ─────────────────────────────────────
\section{Problem Formulation}

\begin{polyamodel}

\textbf{Problem Type}: Find \hfill
\textbf{Domain}: Integer Linear Programming

\textbf{Named Problem}: Knapsack Problem

\end{polyamodel}

% ── Universe ──────────────────────────────────────────────────────────
\subsection{Universe}

\begin{itemize}
  \item $See problem description$
\end{itemize}

% ── Variables ─────────────────────────────────────────────────────────
\subsection{Variables}

\begin{longtable}{lll}
\toprule
\textbf{Symbol} & \textbf{Meaning} & \textbf{Type / Range} \\
\midrule
\endhead
$x_{ij}$ & assignment decision & $\{0, 1\}$ \\
$n$ & number of binary variables = 10 & $\mathbb{{Z}}^+$ \\
\bottomrule
\end{longtable}

% ── Structure ─────────────────────────────────────────────────────────
\subsection{Structure}

Plan a one-day city trip. You have 10 attractions, each with an enjoyment score,
visit duration (minutes), and entry cost (\$). You have 8 hours (480 minutes) and
a \$100 budget. Select the subset of attractions that maximizes total enjoyment
score without exceeding either budget. This is a variant of the 0/1 knapsack
problem with two capacity constraints (time and money).

% ── Mapping ───────────────────────────────────────────────────────────
\subsection{Real-World Mapping}

\begin{longtable}{ll}
\toprule
\textbf{Real-World Concept} & \textbf{Mathematical Object} \\
\midrule
\endhead
problem input & $instance parameters$ \\
\bottomrule
\end{longtable}

% ── Constraints ───────────────────────────────────────────────────────
\subsection{Constraints}

\begin{enumerate}
  \item $Multi-dimensional knapsack$
  \hfill \textit{(0/1 knapsack generalized to two resource constraints)}
  \item $Integer Linear Programming$
  \hfill \textit{(exact optimal solution via Branch \& Bound)}
  \item $Binary selection$
  \hfill \textit{(each attraction is either fully visited or skipped)}
  \item $Dual-budget trade-off$
  \hfill \textit{(balancing time and money simultaneously)}
  \item $Sensitivity analysis$
  \hfill \textit{(how the itinerary changes if time or money budgets shift)}
\end{enumerate}

% ── Objective / Claim ─────────────────────────────────────────────────
\subsection{Objective}

\begin{equation}
\max \sum_{i,j} c_{ij} \, x_{ij}
\end{equation}

% ── Approach ──────────────────────────────────────────────────────────
\subsection{Approach}

\begin{description}
  \item[Suggested approach:] 0/1 Knapsack ILP (PuLP/CBC, Branch \& Bound)
  \item[Complexity class:] 
  \item[Available tools:] 0/1 Knapsack ILP
\end{description}
% ── Phase B: Solution ────────────────────────────────────────────────
\section{Solution}

\begin{polyaresult}

\textbf{Answer}: Solution computed

\textbf{Objective Value}: $97.0$

\textbf{Optimal}: Yes \hfill
\textbf{Feasible}: Yes
\textbf{Algorithm}: 0/1 Knapsack ILP (PuLP/CBC, Branch \& Bound) \quad $$

\textbf{Time}: 0.05123833904508501s

\textbf{Certificate}: Branch-and-bound solved to optimality.

\end{polyaresult}

% ── Solution Details ──────────────────────────────────────────────────
\subsection{Solution Details}


Objective value: z* = 97.0

Method: 0/1 Knapsack with two capacity constraints, formulated as Integer Linear
Programming via PuLP/CBC.

- 10 binary variables (one per attraction: visit or skip)
- 2 capacity constraints (total duration <= 480 min, total cost <= \$100)
- Objective: maximize sum of enjoyment scores for selected attractions

% ── Verification ──────────────────────────────────────────────────────
\subsection{Verification}

\begin{longtable}{lcl}
\toprule
\textbf{Check} & \textbf{Status} & \textbf{Detail} \\
\midrule
\endhead
Feasibility & \passmark & All constraints satisfied \\
Optimality & \passmark & 0/1 Knapsack ILP (PuLP/CBC, Branch \& Bound) \\
\bottomrule
\end{longtable}

% ── Phase C: Interpretation ──────────────────────────────────────────
\section{Interpretation}

% ── The Question & The Answer ─────────────────────────────────────────
\subsection{The Question}
Plan a one-day city trip.

\subsection{The Answer}

\begin{polyainsight}[title={Bottom Line}]
The optimal solution has objective value z* = 97.0.
\end{polyainsight}

% ── What This Means ───────────────────────────────────────────────────
\subsection{What This Means}

- Optimal attraction selection maximizing total enjoyment under 8h and \$100
- Budget utilization bars for both time and money
- Per-attraction efficiency analysis (enjoyment per hour and per dollar)
- Sensitivity analysis across budget variations
- Independent verification of all constraints

0/1 Knapsack with two capacity constraints, formulated as Integer Linear
Programming via PuLP/CBC.

- 10 binary variables (one per attraction: visit or skip)
- 2 capacity constraints (total duration <= 480 min, total cost <= \$100)
- Objective: maximize sum of enjoyment scores for selected attractions

% ── Sensitivity Analysis ──────────────────────────────────────────────

% ── Figures ───────────────────────────────────────────────────────────

% ── Recommendations ───────────────────────────────────────────────────
\subsection{Recommendations}

\begin{enumerate}
  \item The solution is provably optimal; no further improvement is possible within this model.
\end{enumerate}

% ── Limitations ───────────────────────────────────────────────────────
\subsection{Limitations}

\begin{polyawarnbox}
\begin{itemize}
  \item Model assumes deterministic parameters; real-world variability not captured.
  \item Solution quality depends on the accuracy of input data.
\end{itemize}
\end{polyawarnbox}

% ── Appendix ─────────────────────────────────────────────────────────

\end{document}